\infographicchapter{CT Perfusion}{ctp.png}
\label{ch:infographic-ctp}

\section{Overview}

CT perfusion (CTP) is a fast, specialized brain CT scan that measures how well blood is flowing through tiny microscopic vessels inside brain tissue. Instead of taking just one snapshot, the scanner takes a rapid series of low-dose pictures over about 45--60 seconds while a small wave of iodine contrast dye passes through the brain for the first time. By watching how quickly the dye arrives, peaks, and washes out in every part of the brain, smart computer software builds color-coded maps of four key blood-flow numbers: cerebral blood flow (CBF, how fast blood moves through), cerebral blood volume (CBV, how much blood fills the tiny vessels), mean transit time (MTT, how long blood takes to pass through), and Tmax (how delayed the dye's arrival is).

CTP is used most often in emergency stroke care right alongside a standard head CT and CT angiography. Its main job is to tell two areas of the brain apart: tissue that is already permanently injured (the ``infarct core,'' where blood flow has dropped below 30\% of normal and blood volume is very low) versus struggling, low-flow tissue that is still alive and can be rescued (the ``ischemic penumbra,'' where blood arrival is delayed by more than 6 seconds, yet the brain's tiny vessels have widened to keep blood volume up). Comparing those two zones---called core--penumbra mismatch---helps stroke doctors decide whether clot-busting medicine or catheter clot removal (thrombectomy) can still save brain function, even up to 24 hours after symptoms began or when a patient wakes up with stroke symptoms. Doctors also occasionally use it to check for tightened brain arteries after a brain aneurysm bleed (vasospasm) or to study blood supply in brain tumors.

Its biggest strengths are speed right on the emergency CT scanner and fast, automated color maps that guide urgent stroke treatment within minutes. Its main limits are sensitivity to head movement, thick skull shadows, slow contrast flow from weak heart pumping or a tightly narrowed neck artery, and rare overestimation of permanent injury in the very first hours of a stroke. Because taking repeated pictures over the same slice of brain increases local radiation dose and adds extra iodine dye, teams use low-energy settings, strict dose checks, and careful allergy and kidney screening.

\section{Underlying Physical Principles}
\begin{itemize}
    \item \textbf{Watching the First Wave of Dye:} The scanner takes repeated low-dose pictures (usually at 70--80 kVp, which makes iodine contrast pop while keeping local radiation lower) every 1--2 seconds over 45--60 seconds as a compact wave of IV dye sweeps through the brain's tiny vessels.
    \item \textbf{Tracking Brightness Over Time:} Because each spot of brain brightens in direct proportion to how much iodine dye is inside its blood vessels, the series of pictures traces a smooth rise-and-fall curve for every tiny patch of brain tissue.
    \item \textbf{Comparing Against Main Input and Output Vessels:} The software picks a healthy main brain artery (arterial input) and a large draining brain vein (venous output) as references. Comparing brain tissue against those references automatically adjusts for differences in each person's heart rate, heart strength, and injection speed.
    \item \textbf{Calculating True Tissue Blood Flow:} Advanced computer modeling separates the timing of the incoming dye wave from how long blood actually lingers inside each patch of brain tissue, producing accurate blood-flow curves even when blood takes a longer detour around a blockage.
    \item \textbf{Color-Coded Blood-Flow Maps:} From those curves, the computer paints side-by-side color maps of blood flow (CBF), blood volume (CBV), transit time (MTT), and arrival delay (Tmax) so doctors can immediately spot already-injured core versus salvageable brain tissue.
\end{itemize}

\section{Imaging Mechanism}
\begin{enumerate}
    \item \textbf{Gentle Head Support and Target Zone:} Right after a standard head CT, foam pads or a soft forehead strap cradle your head gently so it stays still. The technologist selects an 8--16 cm band of brain tissue to watch and connects a sturdy arm IV to the power injector.
    \item \textbf{Rapid Contrast Wave and 1-Minute Movie:} A compact 40--50 mL wave of iodine dye injects briskly (4--6 mL per second) while the scanner takes low-dose pictures every 1--2 seconds for 45--60 seconds to capture the dye flowing in and washing out.
    \item \textbf{Aligning Every Frame:} The scanner captures 25--40 rapid 3D frames, and automatic software corrects tiny head shifts so every spot of brain stays lined up across the whole minute.
    \item \textbf{Automated Blood-Flow Math:} Perfusion software selects reference brain arteries and veins, traces each voxel's brightness curve, and calculates CBF, CBV, MTT, and Tmax across the brain.
    \item \textbf{Color Maps and Salvageable Tissue Summary:} Within minutes, the workstation shows color blood-flow maps alongside clear volume numbers for permanently injured core (flow below 30\%) and at-risk salvageable tissue (arrival delay over 6 seconds) to guide immediate stroke care.
\end{enumerate}

\section{Key Technical Parameters}
\noindent
\begin{tabularx}{\textwidth}{@{}l X@{}}
\toprule
\textbf{Parameter} & \textbf{Typical Value / Description} \\
\midrule
Physical Principle & Rapid 1-minute series of low-dose CT pictures tracking IV dye flow through brain tissue \\
Energy / Signal Source & Wide-detector CT scanner with a high-flow power injector \\
Typical Frequency / Energy & Low X-ray energy 70--80 kVp (100--200 mAs) to boost iodine brightness and limit local dose \\
Spatial Resolution & About 0.5--1.0 mm detail reconstructed in 5 mm slices covering 8--16 cm of the brain \\
Typical Scan Duration & 45--60 seconds of active scanning (5--15 minutes total for the full stroke CT set) \\
Contrast Agent & Fast IV iodine contrast wave (typically 40--50 mL at 4--6 mL/s followed by saline flush) \\
Ionizing Radiation & Yes (typically 2--5 mSv overall; local head dose carefully tracked per FDA/ACR limits) \\
\bottomrule
\end{tabularx}

\section{Clinical Usages}

\subsection{Primary Usages}
\begin{itemize}
    \item \textbf{Selecting Extended-Window Clot Removal (Thrombectomy):} Measures injured core versus salvageable brain tissue to identify stroke patients who can still benefit from catheter clot removal 6--24 hours after symptoms began or after waking up with a stroke.
    \item \textbf{Guiding Clot-Busting Medicine (Thrombolysis):} Spots favorable salvageable brain tissue in selected patients with late or uncertain symptom start times to guide safe clot-busting IV medicine.
    \item \textbf{Catching Subtle Branch Clots:} Helps doctors spot small, hard-to-see blocked side branches or early low-flow areas that might look nearly normal on a standard head CT or CTA.
    \item \textbf{Checking Natural Backup Circulation:} Shows how well nearby backup vessels (collaterals) are keeping starved brain tissue alive while treatment is being arranged.
\end{itemize}

\subsection{Secondary Usages}
\begin{itemize}
    \item \textbf{Spotting Artery Spasm After Brain Bleeding:} Detects reduced brain blood flow from tightened arteries (vasospasm) days after a ruptured brain aneurysm so doctors can boost blood flow or relax the vessels promptly.
    \item \textbf{Checking Long-Term Artery Narrowing:} Measures blood-flow reserve---sometimes using a safe vessel-widening medicine (acetazolamide)---in chronic carotid blockage or Moyamoya disease to see whether bypass surgery would help.
    \item \textbf{Telling True Stroke Apart From Stroke Look-Alikes:} Helps distinguish a real clot stroke from conditions that mimic stroke, such as temporary weakness after a seizure (Todd paralysis) or a complex migraine.
    \item \textbf{Brain Tumor Blood Supply:} Occasionally checks blood-vessel growth inside brain tumors or helps tell returning tumor apart from radiation scarring when perfusion MRI cannot be done.
\end{itemize}

\section{Advantages}
\begin{itemize}
    \item \textbf{Shows Which Brain Tissue Can Be Saved:} Clear numerical maps (CBF, CBV, MTT, and Tmax) separate permanently injured tissue from living, rescuable tissue to guide life-changing late-window stroke treatments.
    \item \textbf{Built Right Into Emergency Stroke Care:} Runs immediately after standard head CT and CTA on the exact same emergency scanner, adding only a few minutes to stroke evaluation.
    \item \textbf{Fast, Standardized Automated Maps:} Proven commercial software calculates core-and-penumbra volumes within minutes so stroke teams can act quickly and consistently.
    \item \textbf{Catches Subtle Clots:} Picks up small outer-branch blockages and gauges backup blood-flow strength better than plain CT or single-pass CTA alone.
    \item \textbf{Available When MRI Is Not:} Delivers vital blood-flow data 24/7 when urgent brain MRI is unavailable, too slow, or unsafe because of metal implants or unstable vital signs.
\end{itemize}

\section{Disadvantages}
\begin{itemize}
    \item \textbf{Higher Local Radiation Dose:} Taking repeated pictures over the same band of the head gives a higher localized dose than a single head CT; improper repeat scans can risk temporary hair thinning or eye-lens exposure, so dose limits are tightly enforced.
    \item \textbf{Extra Contrast Dye:} Adds another rapid wave of iodine dye on top of CTA, increasing total contrast volume and small risks of dye allergy or kidney stress.
    \item \textbf{Sensitive to Movement and Slow Flow:} Head motion during the 45--60 second scan, weak heart pumping, or a tightly narrowed neck artery can skew reference curves and distort the color maps.
    \item \textbf{Occasional Automated Pitfalls:} Very early after a stroke starts (within 1--3 hours), software can occasionally overestimate permanent injury (a ``ghost infarct core''), and tiny deep strokes or brainstem strokes are harder to see.
    \item \textbf{Smaller Coverage on Older Scanners:} Older scanners with narrower detectors cannot cover the whole brain at once and may miss low-flow spots near the top or very bottom of the brain.
\end{itemize}

\section{Contraindications}

\subsection{Absolute}
\begin{itemize}
    \item \textbf{None in Emergency Stroke Care:} When rapid stroke treatment depends on knowing how much brain tissue can be saved, the urgent benefit of preventing permanent disability generally outweighs standard contrast and radiation concerns.
\end{itemize}

\subsection{Relative}
\begin{itemize}
    \item \textbf{Severe Dye Allergy or Severe Kidney Disease:} A past severe iodine contrast reaction or severe chronic kidney disease (eGFR below 30 mL/min/1.73 m$^2$) favors brain MRI if it is immediately available without delaying emergency stroke care.
    \item \textbf{Pregnancy and Repeat High-Dose Scanning:} Non-urgent CTP is deferred during pregnancy, and repeating CTP over the same brain band is avoided to prevent local skin redness, temporary hair loss, or eye-lens injury.
    \item \textbf{Severe Agitation or Unstable Breathing:} Uncontrollable head movement or an unsafe airway prevents usable blood-flow math and must be stabilized first.
\end{itemize}

\section{Patient Safety Guidelines}
\begin{itemize}
    \item \textbf{Strict Local Radiation Controls:} Because repeated snapshots focus on one band of the head, clinics strictly follow FDA, ACR, and ICRP rules---using low X-ray energy (70--80 kVp), low current, regular dose audits, and avoiding repeat scans---to prevent skin redness, temporary hair thinning, or eye-lens exposure.
    \item \textbf{Balancing Dye Safety With Emergency Stroke Speed:} Staff track total iodine dye when CTP follows CTA and ask about allergies and kidney disease. Per ACR and stroke-society rules, urgent stroke imaging is never delayed waiting for a kidney blood test in patients without known severe kidney disease.
    \item \textbf{Secure High-Flow IV Line:} Checking a sturdy arm IV beforehand prevents dye leakage into arm tissue and avoids cut-off dye waves that would waste radiation.
    \item \textbf{Gentle Head Support and Close Watching:} Soft head immobilization prevents motion blur, while nurses and doctors watch stroke patients closely for breathing trouble, nausea, or sudden symptom changes throughout the scan.
    \item \textbf{Only Used When It Guides Care:} Doctors reserve CTP for cases where blood-flow numbers directly guide treatment (such as late-window stroke care or artery spasm) and check pregnancy status whenever possible without delaying life-saving care.
\end{itemize}

\section{Prerequisite Tests}
\begin{itemize}
    \item \textbf{Noncontrast Head CT and CTA:} In emergency stroke, a plain head CT (to rule out bleeding or a large finished stroke) and head-and-neck CTA (to find a blocked artery) are done right before or alongside CTP as one coordinated package.
    \item \textbf{Quick Stroke Exam and Blood Sugar Check:} A rapid bedside stroke score (NIHSS), a finger-stick blood sugar test (since very low blood sugar can mimic a stroke), and knowing the exact time the person was last seen normal help decide whether CTP is needed.
    \item \textbf{Allergy and Kidney Screening:} Per ACR and AHA/ASA stroke guidelines, teams ask about severe kidney disease and dye allergies, but do not delay urgent stroke CTP waiting for lab creatinine results in patients with no known severe kidney disease.
    \item \textbf{Stable Breathing and Sturdy Arm IV:} Staff confirm a safe airway, stable vital signs, and a secure (18--20 gauge) arm IV line, and check pregnancy status when feasible without slowing emergency treatment.
\end{itemize}

\section{Preparing for the Scan}
\begin{itemize}
    \item \textbf{Emergency vs.\ Planned Visits:} In a stroke emergency, there is zero fasting delay---speed saves brain tissue. For planned outpatient CTP (such as checking long-term blood-flow reserve or a brain tumor), clinics usually ask you not to eat solid food for 3--4 hours while drinking plain water.
    \item \textbf{Medicines and Dye History (Planned Scans):} For elective exams, take regular daily pills with small sips of water, finish any allergy premedication prescribed for a past dye reaction, and tell the team about any recent contrast scans or kidney issues (including metformin guidance if eGFR is below 30 mL/min/1.73 m$^2$).
    \item \textbf{Remove Head and Neck Metal:} Take off glasses, earrings, hairpins, hearing aids, and removable dentures so metal streaks do not cross the brain or reference vessels.
    \item \textbf{Arm IV and Cushioned Head Holder:} A sturdy (18--20 gauge) arm IV delivers the fast dye wave, and your head rests inside a cushioned cradle with soft foam pads or a gentle forehead strap so you can stay motionless effortlessly for 45--60 seconds.
    \item \textbf{Staying Still and Going Home:} Avoid swallowing, talking, or lifting your head during the 1-minute scan, and expect a brief warm flush from the dye. Emergency stroke patients stay under close hospital care; planned outpatients who receive anxiety medicine, drowsy allergy premedication, or an acetazolamide challenge need a responsible adult to drive them home.
\end{itemize}

\section{During Scan}
\begin{itemize}
    \item \textbf{Cushioned Head Support:} You lie on your back with your head cradled in a soft holder and gently steadied with foam pads or a light forehead strap. Staying completely still for just under a minute keeps the blood-flow calculations accurate.
    \item \textbf{Quiet Breathing Without Swallowing:} Breathe softly and steadily without speaking, coughing, or swallowing, because throat and jaw motion can shift your head slightly.
    \item \textbf{What You Will Hear and Feel:} You will hear a steady whirring sound as the table stays still (or glides gently back and forth on wide-coverage scanners) for 45--60 seconds. When the IV dye injects, expect a brief, normal wave of warmth and a slight metallic taste.
    \item \textbf{Close Monitoring and Intercom:} In stroke care, your medical team watches your breathing, comfort, and symptoms closely the entire time. Speak up over the intercom or raise a hand immediately if your IV site hurts or swells, or if you feel nauseous, itchy, or short of breath.
    \item \textbf{How Long It Takes:} The active blood-flow scan takes just 45--60 seconds (usually right after a standard head CT and CTA), with the full combined stroke scan taking 5--15 minutes on the table.
\end{itemize}

\section{Normal Appearance}
\begin{itemize}
    \item \textbf{Balanced Left and Right Sides:} Both halves of the brain look symmetrical across all color blood-flow maps, with no wedge-shaped low-flow zones.
    \item \textbf{Richer Flow in Outer Gray Matter:} The brain's outer gray matter naturally receives more blood flow (about 50--60 mL per 100 g per minute) and holds more blood volume (4--5 mL per 100 g) than the deeper white matter (about 20--25 mL per 100 g per minute and 2 mL per 100 g).
    \item \textbf{Fast, Even Arrival Times:} Transit time (MTT, usually 3--5 seconds) and arrival delay (Tmax, below 4 seconds) stay uniformly short across both sides, with zero injured core or salvageable mismatch volume flagged.
\end{itemize}

\section{Observing for Abnormal Patterns}
\begin{itemize}
    \item \textbf{Already-Injured Tissue (Infarct Core):} Permanently injured brain tissue shows a matching drop in both blood flow (below 30\% of the healthy side) and blood volume, along with severely delayed arrival times.
    \item \textbf{At-Risk Salvageable Tissue (Ischemic Penumbra):} Struggling tissue that can still be saved shows long arrival delays (Tmax over 6 seconds) and moderately lowered blood flow, but normal or even slightly raised blood volume as tiny vessels widen to fight for oxygen---creating a strong mismatch that favors emergency clot treatment.
    \item \textbf{Mildly Slow Flow vs.\ Weak Backup Vessels:} Slightly delayed arrival (Tmax 4--6 seconds) with normal flow and volume is usually harmless tissue that will survive on its own. By contrast, a large share of severely delayed tissue (Tmax over 10 seconds) signals weak backup circulation and faster stroke growth.
    \item \textbf{Extra-High Blood Flow and Stroke Mimics:} Brightly elevated blood flow and volume with fast transit times can appear right after a clot breaks up on its own or during active seizure activity, while post-seizure weakness or complex migraines often show unusual low-flow patches that do not match a single artery's territory.
    \item \textbf{Technical Look-Alikes:} Scanning very early (within 1--3 hours) can occasionally show an exaggerated ``ghost core'' that actually recovers after fast clot removal. Head motion, cut-off dye waves, weak heart pumping, or a tightly narrowed neck artery can also create false wide-area arrival delays.
\end{itemize}

\section{Dealing with Uncertainties}
\begin{itemize}
    \item \textbf{Motion, Cut-Off Waves, and Neck Narrowing:} Head movement during the 1-minute scan, slow dye arrival from weak heart pumping, or a tightly narrowed neck artery can confuse reference curves and make a whole side of the brain look artificially delayed.
    \item \textbf{``Ghost Core'' and Software Differences:} Automated thresholds (like flow below 30\%) can occasionally overestimate permanent damage in the first 1--3 hours after symptoms start, and different software brands can give slightly different numbers---so doctors never reject a patient from life-saving treatment on one automated number alone.
    \item \textbf{Small Deep or Brainstem Strokes:} Very small deep strokes (lacunar infarcts), brainstem strokes, and old scars are hard for CTP to see clearly, and seizures or migraines can create confusing non-arterial patterns.
    \item \textbf{Checking the Raw Data and the Patient:} Stroke teams always look beyond the color summary---checking raw dye curves, motion logs, plain CT slices, CTA vessel pictures, and the patient's actual bedside exam together.
    \item \textbf{Confirming With Brain MRI:} If CTP images are blurry or do not match the patient's symptoms, doctors rely on the bedside exam, plain CT, and CTA to move ahead with urgent care, and use diffusion-weighted brain MRI afterward to confirm the exact final stroke area or rule out stroke look-alikes.
\end{itemize}

\section{Augmenting Methods}
\begin{itemize}
    \item \textbf{Plain Head CT and CTA:} Work together with CTP as the complete emergency stroke triad---plain CT rules out bleeding, CTA pinpoints the blocked artery, and CTP shows how much brain tissue can still be saved.
    \item \textbf{Brain MRI (Diffusion and Perfusion):} The gold standard for confirming true permanent stroke injury---especially for tiny deep strokes, brainstem strokes, very early strokes, and separating true strokes from mimics.
    \item \textbf{Transcranial Doppler Ultrasound:} A gentle, repeatable bedside ultrasound that tracks blood-flow speed inside brain arteries to watch for artery spasm after a brain aneurysm bleed.
    \item \textbf{Blood-Flow Reserve Testing (Acetazolamide, SPECT, or PET):} Uses a safe vessel-widening medicine during CTP or nuclear scans to measure backup blood-flow capacity in chronic carotid blockage or Moyamoya disease before bypass surgery.
    \item \textbf{Automated Stroke Software and Bedside Scores:} Validated perfusion programs combine seamlessly with finger-stick blood sugar, the NIHSS bedside stroke exam, and symptom start time to speed up treatment.
\end{itemize}

\section{Comparison with Alternative Modalities}
\begin{itemize}
    \item \textbf{Versus Brain MRI (DWI/PWI):} MRI is the most accurate test for confirming permanent stroke injury (especially small deep, brainstem, or ultra-early strokes) without X-rays. CTP is much faster, available 24/7 in emergency rooms, and fits seamlessly right after plain CT and CTA.
    \item \textbf{Versus Plain CT and CTA Alone:} Plain CT and CTA are often all that is needed within the first 6 hours of a stroke to rule out bleeding and spot a large clot. CTP adds the tissue-level blood-flow mismatch needed to safely treat strokes in the 6--24 hour extended or wake-up window.
    \item \textbf{Versus Nuclear Scans (SPECT/PET):} PET and SPECT measure absolute brain blood flow and reserve very well in chronic vessel disease, but they are too slow and not available for emergency stroke triage.
\end{itemize}

\section{Follow-up Steps}
\begin{itemize}
    \item \textbf{Immediate Stroke-Team Care and Dye Safety:} Right after the scan, nurses keep a close eye on vital signs and neurologic symptoms, check the arm IV site before removing it, and track total contrast dye with follow-up kidney blood tests over 48--72 hours if you have kidney risk factors.
    \item \textbf{Rapid Treatment Decisions:} The stroke team combines the CTP color maps with plain CT, CTA, your exam score, and symptom timing to start clot-busting IV medicine or move straight to catheter clot removal without delay.
    \item \textbf{Checkup Scans After Treatment:} A follow-up plain head CT (typically 24 hours later, or immediately if symptoms worsen) checks for bleeding or swelling after stroke treatment, and a brain MRI is often done to confirm the final stroke size or clarify stroke mimics.
    \item \textbf{Preventing Future Strokes:} Stroke patients move into specialized stroke-unit care for heart rhythm checks, heart ultrasound, and neck-artery evaluation to prevent future clots. Patients watched for artery spasm after aneurysm bleeding continue bedside Doppler checks and exams.
\end{itemize}

\section{What You Should Know About CT Perfusion?}
\begin{itemize}
    \item \textbf{Painless and Closely Monitored:} You lie on your back with your head gently cushioned and held still for under a minute while IV contrast flows through an arm vein, causing only a brief, normal warm flush and metallic taste.
    \item \textbf{Controlled Radiation and Contrast Safety:} Because CTP takes repeated rapid snapshots of the brain, modern scanners use low-energy X-ray settings and strict dose monitoring; in acute stroke emergencies, the urgent benefit of saving brain tissue far outweighs small contrast or radiation risks.
    \item \textbf{Under One Minute of Scanning:} The active perfusion scan lasts just 45--60 seconds (typically right after a standard head CT and CTA), and automated software generates color-coded brain blood-flow maps within minutes for immediate stroke-team decisions.
    \item \textbf{Why Your Doctor Orders It:} In suspected stroke, CTP shows not only which brain tissue is already injured, but also which low-flow brain tissue can still be rescued by clot-busting medicine or catheter clot removal.
    \item \textbf{What to Tell Your Care Team:} In an emergency, sharing the exact time symptoms started (or when you were last seen well) is critical; for elective scans, tell the team about any contrast allergy, kidney disease, or recent contrast scans, and report any IV site pain or breathing trouble immediately.
\end{itemize}

\section{Key References}
\begin{itemize}
    \item Powers WJ, Rabinstein AA, Ackerson T, et al. \textit{Guidelines for the Early Management of Patients With Acute Ischemic Stroke: An Update to the AHA/ASA Guidelines}. Stroke.
    \item American College of Radiology (ACR). \textit{ACR Appropriateness Criteria: Cerebrovascular Diseases---Stroke and Stroke-Related Conditions} and \textit{ACR--ASNR--SPR Practice Parameter for the Performance of Computed Tomography (CT) Perfusion in Neuroradiologic Imaging}.
    \item U.S. Food and Drug Administration (FDA). \textit{Safety Investigation of CT Brain Perfusion Scans: Radiation Dose Controls}.
    \item StatPearls Publishing. \textit{CT Perfusion in Acute Ischemic Stroke}.
\end{itemize}
